\documentclass[11pt]{article}
\usepackage{explorer}
\newcommand{\ExplorerLanguage}{FR}
\newcommand{\ExplorerNumber}{11}
\newcommand{\ExplorerKind}{Leçon et méthode}
\begin{document}
\ExplorerTitle{Lycée · optimisation et olympiades}{Inégalités et cas d’égalité}
\ExplorerSection{Avant de commencer}
Carrés, factorisation et racines carrées.\par
\ExplorerSection{Idées et méthode}
Pour tout réel \(u, u^{2}\geq 0\), avec égalité exactement lorsque \(u=0\). Beaucoup d’inégalités se ramènent à cette observation.\par
Pour \(a,b\geq 0, (\sqrt{a}-\sqrt{b})^{2}\geq 0\) donne \(a+b\geq 2\sqrt{ab}\). C’est l’inégalité entre les moyennes arithmétique et géométrique. L’égalité a lieu si et seulement si \(a=b\).\par
Prouver une borne ne suffit pas à annoncer un maximum : il faut vérifier qu’un cas d’égalité autorisé l’atteint. Les contraintes de positivité et de domaine font partie du raisonnement.\par
\ExplorerSection{Exemple expliqué}
Si \(a+b=10\) avec \(a,b\geq 0\), alors \((a-b)^{2}\geq 0\) donne \(100-4ab\geq 0\), donc \(ab\leq 25\). Le couple \(a=b=5\) atteint 25 : c’est bien le maximum.\par
\ExplorerSection{À explorer}
Tracer \(f(x)=x(10-x)\) pour \(0\leq x\leq 10\) et déplacer le point. Retrouver ensuite le maximum sans utiliser le dessin.\par
\end{document}
